\chapter{The Agentic Coding Workflow}
\label{ch:agentic}

The change this chapter describes is not that code gets written faster. It is
that the job changes. You stop being the person who writes the code and become
the person who specifies it, constrains the environment it runs in, and verifies
what came out. That is closer to system administration than to programming, and
the skills that transfer are the operational ones: permissions, logging,
reproducibility, knowing what you would do if it went wrong.

\section{Terminology, briefly}

\term{Vibe coding} was coined in early 2025 for accepting AI-generated code
without close review --- going with the flow. It has since been diluted into
shorthand for any AI-assisted coding, which loses the distinction.

\term{Chat-oriented programming (CHOP)} is the more useful term: an iterative
loop in which you describe intent, the system proposes code, you run it, and you
refine. The distinguishing feature is that you read and verify at each step.

The difference is whether verification is in the loop. Vibe coding is fine for a
throwaway plot. It is not fine for anything entering a manuscript, a database,
or a pipeline someone else will run.

\begin{keypoint}
The question is not whether AI wrote the code. It is whether anyone checked it.
You are accountable for what you publish regardless of what produced it.
\end{keypoint}

\section{Architecture: chatbot versus harness}

\begin{table}[htbp]
\centering
\caption{Two interaction models, two risk profiles.}
\label{tab:chat-vs-harness}
\small
\begin{tabularx}{\textwidth}{lXX}
\toprule
 & \textbf{Chatbot} & \textbf{Harness} \\
\midrule
Interface     & Web or app & Terminal, in your project \\
Sees          & What you paste & Your files, directly \\
Acts          & No & Runs commands, edits files \\
Verification  & You copy, paste, run & It runs, reads the error, retries \\
Failure cost  & A wrong snippet & Modified or deleted files \\
Good for      & Explaining, drafting & Multi-file work, refactors, pipelines \\
\bottomrule
\end{tabularx}
\end{table}

The capability difference comes from the feedback loop. A harness sees the
traceback. A chatbot sees what you chose to paste, which is usually the last
five lines, which is usually not where the problem is.

\section{The while loop}

\textcite{liu2026claudecode} analysed a production coding agent's source and
found the core to be exactly the loop from \Cref{ch:landscape}: call the model,
execute requested tools, feed results back, repeat until done. Their analysis
traces five motivating values --- human decision authority, safety and security,
reliable execution, capability amplification, contextual adaptability --- through
thirteen design principles to implementation choices.

The distribution of code is the finding worth internalising. Almost none of it is
the loop. The bulk is the permission system, context management, error recovery
and memory. When you evaluate a harness, evaluate those, because that is what
differs.

\section{A file is all you need}

\textcite{piskala2026files} makes the Unix argument for agent design: collapse
heterogeneous resources into uniform file-like interfaces. In practice this means
your project communicates with an agent through files, and those files are
readable by you.

\subsection{README.md}

Written for a human, read by both. It should state what the project does, how to
run it, what the data is and where it lives, what the dependencies are, and
what the current state is. If a new postdoc could not start from it, an agent
cannot either.

\subsection{AGENTS.md}

This is the behavioural contract. Not aspirations --- constraints, written in
response to failures you have observed.

\begin{lstlisting}[caption={An AGENTS.md that earns its place}]
# Conventions

## Environment
- R 4.4 via renv. Activate with `renv::restore()`.
- Never install packages globally. Add to renv.lock.

## Data
- data/raw/ is READ-ONLY. Never write there, never modify.
- All derived data goes to data/processed/ and must be
  regenerable by running scripts/ in numeric order.
- Coordinates are WGS84 decimal degrees. Do not reproject
  silently; if a reprojection is needed, say so first.

## Code
- Tidyverse style. snake_case. No magic numbers.
- Every script starts by reading from data/ and ends by
  writing to results/. No side effects elsewhere.
- Taxon names go through taxonomic backbone matching before
  any join. Never join on raw name strings.

## Don't
- Don't commit. I commit.
- Don't touch manuscript/ unless explicitly asked.
- Don't add dependencies without asking.
\end{lstlisting}

Each line there closes a specific failure. The taxon name rule exists because
joining on unmatched names silently drops records. The reprojection rule exists
because a silent CRS change produces a map that looks fine and is wrong.

\subsection{Skills and packaged instructions}

Several harnesses support bundling instructions, scripts and reference material
into reusable units --- variously called skills, plugins or commands
\parencite{anthropic2026skills}. The useful ones are, again, just directories of
files: a description of when to use them, the procedure, and any scripts needed.
For a laboratory with a standard analysis protocol, packaging it this way is a
reasonable way to make the protocol executable rather than merely documented.

\subsection{Memory}

Some harnesses maintain persistent memory across sessions. When they do, prefer
the implementations that store it as files you can read and edit. Opaque memory
is a reproducibility problem: if the agent behaves differently today and you
cannot see why, you have lost control of your own environment.

\section{Tokens, context and cost}

Tokens are the unit of both billing and context. The approximations from
\Cref{ch:powers}: one token is about four characters or three-quarters of a
word; a 200-line R script is roughly 1,000--2,000 tokens.

Everything in the context window is re-sent on every turn. A long session with a
large file loaded is paying for that file repeatedly. Subscription plans meter
this: a typical consumer tier allocates on the order of tens of thousands of
tokens per rolling five-hour window, which a single large refactor can consume.

Practical economy:

\begin{itemize}
  \item Start a new session per task. Do not carry an hour of unrelated history.
  \item Point at files rather than pasting them; let the agent read what it
    needs.
  \item Ask for a diff, not a rewritten file.
  \item Use a smaller model for mechanical work. Not everything needs the
    frontier model.
  \item Keep a scratch file for intermediate results instead of keeping them in
    context.
\end{itemize}

\section{Prompting, practically}

Prompt engineering as a discipline is overstated for frontier models and
genuinely important for smaller local ones. What consistently helps:

\textbf{State the goal, the constraints, and the definition of done.} ``Fix the
script'' produces guesswork. ``The script fails at line 40 with this traceback;
it should output one row per station with mean abundance; do not change the input
format'' produces a fix.

\textbf{Provide an example.} One worked input--output pair beats three
paragraphs of specification, particularly for data transformation.

\textbf{Ask for a plan first on anything non-trivial.} Reviewing a plan costs a
minute. Reviewing a wrong 400-line implementation costs an hour.

\textbf{Say what not to do.} Negative constraints are effective and
underused. ``Do not modify data/raw. Do not add dependencies. Do not reformat
unrelated code.''

\textbf{Give it a way to check itself.} A test, an assertion, an expected row
count. An agent that can verify will iterate to correctness; one that cannot will
report success.

\section{The review discipline}

The failure mode that matters is not bad code. It is plausible code that runs,
produces output of the right shape, and is subtly wrong.

\begin{enumerate}
  \item \textbf{Commit before.} Always. A clean checkpoint makes everything
    reversible.
  \item \textbf{Read the diff.} All of it. If it is too large to read, the task
    was too large --- revert and split it.
  \item \textbf{Check the boundaries.} Empty inputs, single rows, missing values,
    duplicate keys. These are where generated code fails.
  \item \textbf{Verify the joins.} Row counts before and after. Silent row loss
    in a merge is the single most common serious error in ecological data
    processing, and agents produce it readily.
  \item \textbf{Run it twice.} Same input, same output. If not, there is hidden
    state or unseeded randomness.
  \item \textbf{Sanity-check the science.} Does the magnitude make sense? Are
    the units right? Would you have expected this? The agent has no domain
    intuition and will not flag a biologically impossible result.
\end{enumerate}

\begin{keypoint}
Agents fail loudly on syntax and quietly on semantics. Your review should spend
its time on the second kind.
\end{keypoint}

\section{Scope control}

The most reliable determinant of success is task size. Agents are good at
bounded, well-specified tasks and degrade sharply as scope grows --- which is
what \textcite{guo2026promptbio} measured across 244 bioinformatics tasks.

Good scope: ``write a function that reads this CSV and returns a tidy data frame
with these columns.'' Bad scope: ``build the analysis pipeline.'' The second is a
sequence of the first, and running it as a sequence gives you a review point
between each step.
